\begin{center}
{\Large \bfseries Optimization of storage for flood simulation data using sparse structures}

\vspace{1cm}
{\Large \bfseries ABSTRACT}

\vspace{1.5cm}
\end{center}


River floods are among the most costly natural disasters, and their study increasingly relies on high-resolution hydraulic simulations that predict how water spreads across the terrain. These simulations generate such large volumes of data that their subsequent analysis becomes a problem in itself: answering questions such as when water reaches a given area or which streets become dangerous for people requires scanning hundreds of gigabytes of raster files, with waits of minutes or hours. However, at any instant of a flood the vast majority of the territory remains dry, and storing all those empty cells is a waste that this work turns into an opportunity.

With this motivation, this work designs and implements an efficient storage and analysis system for high-resolution hydraulic simulations. Simulations from the Triton simulator generate up to 235 GB of GeoTIFF data per scenario; the developed ETL pipeline transforms them into a TileDB sparse array with a ByteShuffle+ZSTD-9 filter, reducing them to 17.2 GB (a 92.7\,\% reduction, 13.7$\times$ ratio) by storing only the wet cells, which account for just $\sim$8.3\,\% of the domain.

The system includes an engine of 24 analytical queries organized in eight blocks, ranging from basic water-depth and velocity queries to indicators of hydrodynamic instability, the severe-damage zone (Spanish Royal Decree 9/2008) and the available evacuation time. A benchmark against GeoTIFF shows advantages of up to 12.4$\times$ in five of six access patterns, consistent across four datasets and two different machines. Functional validation confirms the accuracy of the queries against the original data.

The system has been applied to 56 simulations over a basin of $\sim$78,000 km$^2$ in the United States, and is generalizable to any equivalent 2D hydraulic simulator. The interface is a Streamlit web application with interactive maps, multi-scenario comparison and export to standard GIS formats (GeoTIFF and QGIS), deployed in production.


